%%%%%%%%%%%%%%%%%%%%%%%%%%%%%%%%%%%%%%%%%%%%%%%%%%%%%%%%%%%%%
% 通用附录模板 (v1.5.5 重写, 通用化) — 替代原 2024B 烟幕题 example
%
% 设计原则:
% 1. 整段都是 \TeX 注释 (以 % 开头), check 16 跳过 % 行 → dryrun 不报 RED
% 2. 留 1 行 INLINE 占位符 \AppendixSlot{...}, check 16 抓 \texttt{} 引用 → 跑题用户填完才 GREEN
% 3. 跑题用户复制本文件后, 删注释 + 填实际文件路径, 删 \AppendixSlot 行
%
% 2026 规范第五/十一条 + 2026-08-26 官方答疑 Q3:
% - 附录首页应是支撑材料清单 (file list)
% - 附录 1: 程序使用说明 (用了/没用)
% - 附录 2: 支撑材料说明 (提供了/没有)
% - 附录 3: 核心源代码 (推荐)
% - ⚠️ 漏写就是 0 分 (评审硬性要求)
%%%%%%%%%%%%%%%%%%%%%%%%%%%%%%%%%%%%%%%%%%%%%%%%%%%%%%%%%%%%%

\section*{附录 1：程序使用说明}

%%\subsection*{运行环境}
%%本参赛队使用 Python \texttt{3.12} + 主要依赖: numpy $\ge 1.24$, scipy $\ge 1.11$, pandas $\ge 2.0$, matplotlib $\ge 3.7$, openpyxl $\ge 3.1$。如果用了其他库 (如 scikit-learn, PuLP, PyMC), 在此列出。

%%\subsection*{目录结构}
%%\texttt{求解/} 目录下的主要文件:
%%\begin{itemize}
%%  \item \texttt{共享/constants.py} --- 跨问题物理常量
%%  \item \texttt{共享/params.py} --- 跨问题复用参数
%%  \item \texttt{问题1/问题1\_xxx.py} --- 问题 1 主代码
%%  \item \texttt{问题1/图片/*.png} --- 问题 1 输出图
%%  \item \texttt{问题1/结果/*.csv} --- 问题 1 输出表
%%  ... (按你实际问题列, 1-2 行/文件)
%%\end{itemize}

%%\subsection*{运行方法}
%%\begin{verbatim}
%%cd 求解
%%python 问题1/问题1_xxx.py
%%python 问题2/问题2_yyy.py
%%... (按你实际问题列)
%%\end{verbatim}

\section*{附录 2：支撑材料说明}

%%支撑材料包 \texttt{支撑材料.rar} 内容:
%%\begin{enumerate}
%%  \item \texttt{求解/}: 所有 Python 源代码
%%  \item \texttt{求解/问题*/图片/}: 所有图表 (PNG)
%%  \item \texttt{求解/问题*/结果/}: 所有结果 (CSV/XLSX)
%%  \item \texttt{AI工具使用详情.pdf}: AI 工具使用情况详细说明
%%  \item (如果用) \texttt{原始数据.xlsx}: 题面附件
%%\end{enumerate}
%%支撑材料大小: 约 X MB (远小于 20 MB 上限)。

\section*{附录 3：核心源代码}

%% 官方 Word 模板建议: 附录 3 放核心源代码 (贴关键文件 / 节选片段),
%% 配合"该代码使用 AI 工具: ..."标注。挑 1-2 个最关键函数贴入,
%% 不用整文件贴 (会超页)。

%%\subsection*{关键函数 1 (xxx.py 第 N 行)}
%%该部分代码使用 AI 工具: 【工具名, 版本, 日期】 / 团队手写
%%\begin{verbatim}
%%def your_key_function(...):
%%    """功能描述"""
%%    ... (贴核心逻辑, 30-50 行内)
%%\end{verbatim}

% ===== 跑题用户最终要填的实际内容 =====
% 把上面所有 % 注释替换为你的实际内容, 然后在下面填 1 行 inline 占位符
% 让 dryrun check 16 验证你的支撑材料文件清单是否都在

\textbf{【在这里列你实际的 .py / .png / .xlsx 文件路径, 跟支撑材料.zip 完全对齐。dryrun check 16 会逐个 Test-Path 验证。本行留【】= dryrun check 15 RED, 必须替换】}
